\chapter{Эрмитовы наблюдаемые, спин и правило вероятности}

\lectureoverview{Квантовая механика — это исчисление вероятностей. Чтобы пользоваться ею, не подменяя понимание лозунгами, нам нужна компактная математическая цепочка: комплексные фазы, нормированные состояния, эрмитовы наблюдаемые, собственные значения и собственные векторы, ортогональность и правило Борна. Затем три матрицы Паули превращают эту цепочку в конкретные предсказания для спина электрона.}

\section{Комплексные числа в полярной форме}

Начнём с устранения одного небольшого источника алгебраических затруднений. Комплексное число
\begin{equation}
  z=x+iy
\end{equation}
— это точка на комплексной плоскости, где \(x\) лежит на действительной оси, а \(y\) — на мнимой. Если эта точка имеет полярные координаты \((r,\theta)\), то
\begin{equation}
  x=r\cos\theta,
  \qquad
  y=r\sin\theta,
\end{equation}
и, следовательно,
\begin{equation}
  z=r(\cos\theta+i\sin\theta)=re^{i\theta}.
  \label{eq:l03-polar-form}
\end{equation}
Обозначение Эйлера полезно потому, что формулы сложения углов принимают вид показательного закона
\begin{equation}
  e^{i\theta}e^{i\phi}
  =e^{i(\theta+\phi)}.
  \label{eq:l03-phase-product}
\end{equation}
В самом деле, перемножая две тригонометрические формы, получаем
\begin{align}
& (\cos\theta+i\sin\theta)(\cos\phi+i\sin\phi)
\notag\\
&\qquad=\cos(\theta+\phi)+i\sin(\theta+\phi).
\end{align}
Все привычные формулы суммы содержатся в этом одном умножении.

Комплексное сопряжение отражает точку относительно действительной оси:
\begin{equation}
  z^*=x-iy,
  \qquad
  (e^{i\theta})^*=e^{-i\theta}.
\end{equation}
Следовательно,
\begin{equation}
  e^{i\theta}e^{-i\theta}=1.
\end{equation}
Числа вида \(e^{i\theta}\) лежат на единичной окружности и называются \emph{чистыми фазами}. Любое ненулевое комплексное число равно положительному модулю, умноженному на чистую фазу:
\begin{equation}
  z=|z|e^{i\theta},
  \qquad
  |z|^2=z^*z=x^2+y^2.
  \label{eq:l03-modulus-phase}
\end{equation}

\begin{classroomqa}
\lecturetimestamp{00:06:36}
\audiencequestion{Соответствует ли замена \(\theta\) на \(-\theta\) обращению направления вращения на комплексной плоскости?}
\lecturerresponse{Да. Положительные и отрицательные фазовые углы обходят единичную окружность в противоположных направлениях. Комплексное сопряжение превращает \(e^{i\theta}\) в \(e^{-i\theta}\), что и есть отражённая точка и обратная фаза.}
\end{classroomqa}

На сегодня из комплексных чисел нам нужно только это: модуль, фаза, сопряжение и единичная окружность.

\section{Квантовая механика как исчисление вероятностей}

Теперь перейдём к постулатам. Физическую систему подготавливают в \emph{состоянии}; измерение запрашивает значение \emph{наблюдаемой}; квантовая механика задаёт вероятности возможных ответов.

Состояния — это векторы в комплексном векторном пространстве. Для двухуровневой спиновой системы
\begin{equation}
  \ket\psi=\alpha\ket{u}+\beta\ket{d},
  \qquad
  \ket{u}=\begin{pmatrix}1\\0\end{pmatrix},
  \qquad
  \ket{d}=\begin{pmatrix}0\\1\end{pmatrix}.
  \label{eq:l03-two-level-state}
\end{equation}
Коэффициенты определяют вероятности состояний вверх и вниз,
\begin{equation}
  P(u)=|\alpha|^2,
  \qquad
  P(d)=|\beta|^2,
\end{equation}
поэтому состояние должно быть нормировано:
\begin{equation}
  \braket{\psi}{\psi}
  =|\alpha|^2+|\beta|^2
  =1.
  \label{eq:l03-state-normalization}
\end{equation}
Размерность относится к числу независимых альтернатив в пространстве состояний, а не к размерности обычного пространства. Для системы с тремя взаимно исключающимися базисными состояниями использовались бы три ортонормированных столбца и три амплитуды.

Вероятности не определяют амплитуды полностью. Умножение на чистую фазу не изменяет модуль:
\begin{equation}
  |e^{i\theta}\alpha|^2=|\alpha|^2.
\end{equation}
Следовательно, амплитуды содержат информацию, выходящую за рамки вероятностей в этом одном базисе. Её роль станет видна, когда мы рассмотрим другие измерения и, позже, временную эволюцию.

Базисные состояния ортонормированы:
\begin{equation}
  \braket{u}{u}=\braket{d}{d}=1,
  \qquad
  \braket{u}{d}=\braket{d}{u}=0.
\end{equation}
Их ортогональность относится к комплексному пространству состояний. Пока её не следует путать с двумя стрелками, образующими прямой угол в обычном трёхмерном пространстве; это различие станет решающим ближе к концу лекции.

\section{Почему наблюдаемые эрмитовы}

Экспериментальные результаты можно записывать вещественными числами. Для удобства мы можем упаковать два вещественных числа в одно комплексное число, но сами показания стрелки прибора всегда можно разложить на вещественные величины. В конечномерном случае квантовые наблюдаемые представляются специальным классом линейных операторов: эрмитовыми матрицами.

Чтобы определить их, сначала транспонируем матрицу, меняя местами строки и столбцы:
\begin{equation}
  (M^T)_{ij}=M_{ji}.
  \label{eq:l03-transpose}
\end{equation}
Вещественная матрица, удовлетворяющая условию \(M^T=M\), симметрична. Для комплексных матриц соответствующей операцией служит эрмитово сопряжение,
\begin{equation}
  M^\dagger=(M^T)^*,
  \qquad
  (M^\dagger)_{ij}=M_{ji}^*.
  \label{eq:l03-dagger}
\end{equation}

\begin{figure}[htbp]
\centering
\includegraphics[width=0.88\linewidth]{lecture_03_figure_02.png}
\caption{Транспонирование, симметрия и переход к эрмитову сопряжению располагаются на доске по одной и той же схеме.}
\lectureframe{00:23:06}
\end{figure}

\begin{definition}
Матрица \(M\) называется \emph{эрмитовой}, если
\begin{equation}
  M^\dagger=M,
  \qquad\text{equivalently}\qquad
  M_{ij}=M_{ji}^*.
  \label{eq:l03-hermitian}
\end{equation}
\end{definition}

Диагональные элементы вещественны, поскольку \(M_{ii}=M_{ii}^*\). Внедиагональные элементы, симметричные относительно диагонали, образуют сопряжённые пары. Поэтому общий вид матрицы два на два таков:
\begin{equation}
  M=
  \begin{pmatrix}
    a&z\\
    z^*&b
  \end{pmatrix},
  \qquad a,b\in\mathbb R.
  \label{eq:l03-hermitian-two}
\end{equation}
Эрмитовость — матричный аналог условия вещественности; она не требует, чтобы каждый элемент по отдельности был вещественным.

Первое полезное следствие вытекает из двух элементарных тождеств сопряжения:
\begin{equation}
  \braket{B}{A}=\braket{A}{B}^*,
  \label{eq:l03-inner-conjugation}
\end{equation}
и, для любого \(M\),
\begin{equation}
  \bra B M\ket A
  =
  \bigl(\bra A M^\dagger\ket B\bigr)^*.
  \label{eq:l03-matrix-element-conjugation}
\end{equation}
Выражение \(\bra B M\ket A\) — это число, матричный элемент \(BA\) матрицы \(M\). Если матрица \(M\) эрмитова и \(B=A\), то
\begin{equation}
  \bra A M\ket A
  =
  \bigl(\bra A M\ket A\bigr)^*.
\end{equation}
Следовательно, оно вещественно.

\begin{proposition}
Для любой эрмитовой матрицы \(M\) среднее значение
\begin{equation}
  \expect{M}_A=\bra A M\ket A
  \label{eq:l03-expectation}
\end{equation}
вещественно в любом нормированном состоянии \(\ket A\).
\end{proposition}

Вероятностный смысл среднего значения будет раскрыт далее, но уже его вещественность объясняет, почему эрмитовы операторы являются естественными кандидатами на роль наблюдаемых.

\begin{classroomqa}
\lecturetimestamp{00:33:22}
\audiencequestion{Происходит ли обозначение среднего угловыми скобками из обозначения bra--ket?}
\lecturerresponse{Да. Скобочная нотация Дирака использует те же самые скобки. В \(\bra A M\ket A\) оператор зажат между bra и ket, представляющими состояние, и результатом является среднее значение.}
\end{classroomqa}

\begin{classroomqa}
\lecturetimestamp{00:34:31}
\audiencequestion{Может ли сама физическая величина иметь мнимое измеренное значение? Если результаты измерений вещественны, зачем в квантовой механике нужны комплексные числа?}
\lecturerresponse{Набор измеренных результатов всегда можно записать вещественными числами, даже если мы решим упаковать пару в одно комплексное число. Более точный вопрос состоит в том, почему теория вынуждена использовать комплексные амплитуды. Пока удержите этот вопрос: ответ связан с обратимой временной эволюцией, которая ещё не введена. Если ограничить теорию вещественными векторами, будут потеряны важные возможности.}
\end{classroomqa}

Математика — не украшение. Без неё запутанность можно заставить звучать таинственно, оставляя её при этом неопределённой. Мы продолжаем в максимально простой постановке, с одной двухуровневой системой, чтобы позже можно было честно объединить две такие системы.

\section{Собственные значения — возможные результаты измерения}

\begin{definition}
Ненулевой вектор \(\ket a\) является собственным вектором \(M\) с собственным значением \(\lambda_a\), если
\begin{equation}
  M\ket a=\lambda_a\ket a.
  \label{eq:l03-eigenvalue-equation}
\end{equation}
\end{definition}
Матрица может изменить длину вектора или его фазу, но не превращает его в другое направление в пространстве состояний.

Для эрмитовой матрицы собственное значение вещественно. Умножая уравнение ~\eqref{eq:l03-eigenvalue-equation} слева на \(\bra a\), получаем
\begin{equation}
  \bra a M\ket a
  =\lambda_a\braket{a}{a}.
\end{equation}
Левая часть вещественна в силу эрмитовости, а норма справа вещественна и положительна. Следовательно,
\begin{equation}
  \lambda_a
  =
  \frac{\bra a M\ket a}{\braket{a}{a}}
  \in\mathbb R.
  \label{eq:l03-real-eigenvalue}
\end{equation}

Теперь следует физическое утверждение: когда \(M\) представляет наблюдаемую, соответствующие собственные значения — это значения, которые может дать измерение \(M\). Собственный вектор, соответствующий \(\lambda_a\), есть состояние, в котором это значение получается наверняка.

Простейший пример — диагональная матрица:
\begin{equation}
  M=
  \begin{pmatrix}
    m_{11}&0\\
    0&m_{22}
  \end{pmatrix}.
\end{equation}

\begin{figure}[htbp]
\centering
\includegraphics[width=0.72\linewidth]{lecture_03_figure_03.png}
\caption{Абстрактное уравнение на собственные значения помещено над первым диагональным примером.}
\lectureframe{00:44:04}
\end{figure}

Прямое умножение даёт
\begin{align}
  M\begin{pmatrix}1\\0\end{pmatrix}
  &=m_{11}\begin{pmatrix}1\\0\end{pmatrix},
  \notag\\
  M\begin{pmatrix}0\\1\end{pmatrix}
  &=m_{22}\begin{pmatrix}0\\1\end{pmatrix}.
  \label{eq:l03-diagonal-eigenvectors}
\end{align}

\begin{figure}[htbp]
\centering
\includegraphics[width=0.72\linewidth]{lecture_03_figure_04.png}
\caption{Нижний базисный вектор получается умноженным на нижний диагональный элемент.}
\lectureframe{00:45:58}
\end{figure}

Для диагональной матрицы большего размера действует то же правило: её диагональные элементы являются собственными значениями, а стандартные столбцы с одной единицей — собственными векторами.

\section{Три спиновые матрицы}

Для спина электрона вдоль третьей оси возьмём
\begin{equation}
  \sigma_3=
  \begin{pmatrix}
    1&0\\
    0&-1
  \end{pmatrix}.
\end{equation}
Её собственные состояния — \(\ket u=(1,0)^T\) с собственным значением \(+1\) и \(\ket d=(0,1)^T\) с собственным значением \(-1\). Поэтому двумя возможными показаниями детектора являются эти два собственных значения.

Спиновая матрица для первой оси имеет вид
\begin{equation}
  \sigma_1=
  \begin{pmatrix}
    0&1\\
    1&0
  \end{pmatrix}.
\end{equation}
Её предполагаемые собственные векторы сначала записываются без нормировки, поскольку умножение собственного вектора на любой ненулевой скаляр оставляет его собственным вектором. Однако физические состояния должны быть нормированы:
\begin{align}
  \ket{x+}
  &=\frac{1}{\sqrt2}\begin{pmatrix}1\\1\end{pmatrix},
  &
  \sigma_1\ket{x+}&=+\ket{x+},
  \notag\\
  \ket{x-}
  &=\frac{1}{\sqrt2}\begin{pmatrix}1\\-1\end{pmatrix},
  &
  \sigma_1\ket{x-}&=-\ket{x-}.
  \label{eq:l03-sigma-one-eigenvectors}
\end{align}

\begin{figure}[htbp]
\centering
\includegraphics[width=0.72\linewidth]{lecture_03_figure_05.png}
\caption{Внедиагональная матрица и кандидат с одинаковыми весами появляются ещё до завершения перемножения.}
\lectureframe{00:50:29}
\end{figure}

\begin{classroomqa}
\lecturetimestamp{00:56:45}
\audiencequestion{Можно ли увидеть собственные векторы \(\sigma_1\) геометрически, а не находить их алгебраически?}
\lecturerresponse{В вещественной двумерной картине \(\sigma_1\) меняет местами базисные направления \(E_1\) и \(E_2\). Диагональное направление \((1,1)^T\) остаётся неизменным, тогда как другая диагональ \((1,-1)^T\) меняет знак. Следовательно, это собственные векторы с собственными значениями \(+1\) и \(-1\). Это картина в пространстве компонент, а не пока ещё утверждение о физическом указателе спина.}
\end{classroomqa}

Оставшаяся матрица имеет вид
\begin{equation}
  \sigma_2=
  \begin{pmatrix}
    0&-i\\
    i&0
  \end{pmatrix},
\end{equation}
с нормированными собственными векторами
\begin{equation}
  \ket{y+}=\frac{1}{\sqrt2}\begin{pmatrix}1\\i\end{pmatrix},
  \qquad
  \ket{y-}=\frac{1}{\sqrt2}\begin{pmatrix}1\\-i\end{pmatrix}.
  \label{eq:l03-sigma-two-eigenvectors}
\end{equation}
Их собственные значения снова равны \(+1\) и \(-1\).

Квадрат каждой из трёх матриц равен единичной матрице:
\begin{equation}
  \sigma_1^2=\sigma_2^2=\sigma_3^2=I.
  \label{eq:l03-pauli-squares}
\end{equation}
Если \(\sigma_i\ket a=\lambda\ket a\), то двукратное применение \(\sigma_i\) даёт \(\lambda^2=1\), так что \(\lambda=\pm1\). Совпадение спектров — первый признак того, что эти матрицы представляют компоненты спина вдоль эквивалентных пространственных осей.

\section{Различные результаты требуют ортогональных состояний}

Прежде чем доказывать следующую теорему, сформулируем её смысл. Если одна наблюдаемая даёт в двух состояниях два разных определённых результата, то измерение этой наблюдаемой может идеально различить эти состояния. В квантовой механике такая идеальная различимость выражается ортогональностью.

\begin{theorem}
Пусть \(M\) эрмитов, причём
\begin{equation}
  M\ket A=\lambda_A\ket A,
  \qquad
  M\ket B=\lambda_B\ket B.
\end{equation}
Если \(\lambda_A\ne\lambda_B\), то
\begin{equation}
  \braket{B}{A}=0.
\end{equation}
\end{theorem}

\begin{proof}
Домножение первого уравнения для собственного значения слева на \(\bra B\) даёт
\begin{equation}
  \bra B M\ket A=\lambda_A\braket{B}{A}.
\end{equation}
Эрмитовость и второе уравнение для собственного значения дают тот же матричный элемент в виде
\begin{equation}
  \bra B M\ket A=\lambda_B\braket{B}{A},
\end{equation}
поскольку \(\lambda_B\) вещественно. Вычитая, получаем
\begin{equation}
  (\lambda_A-\lambda_B)\braket{B}{A}=0.
\end{equation}
Собственные значения различны, следовательно, скалярное произведение должно обращаться в нуль.
\end{proof}

Собственные векторы \(\sigma_1\) дают непосредственную проверку:
\begin{equation}
  \braket{x+}{x-}
  =
  \frac12
  \begin{pmatrix}1&1\end{pmatrix}
  \begin{pmatrix}1\\-1\end{pmatrix}
  =\frac12-\frac12=0.
\end{equation}
И наоборот, для двух ортогональных состояний существует измерение, которое их различает. Следовательно, ортогональность — не просто геометрическое слово; она имеет операциональный смысл.

\section{Постулат о вероятности}

Предположим, что система приготовлена в нормированном состоянии \(\ket B\). Пусть \(M\) — наблюдаемая, а \(\ket A\) — нормированный собственный вектор \(M\) с собственным значением \(\lambda_A\). Тогда вероятность того, что измерение \(M\) даст \(\lambda_A\), равна
\begin{equation}
  P(\lambda_A\mid B)
  =
  |\braket{A}{B}|^2
  =
  \braket{A}{B}\braket{B}{A}.
  \label{eq:l03-born-rule}
\end{equation}
Это произведение вещественно и неотрицательно. Для нормированных векторов оно не превосходит единицы.

\begin{figure}[htbp]
\centering
\includegraphics[width=0.78\linewidth]{lecture_03_figure_06.png}
\caption{Правило вероятности записывается как перекрытие, умноженное на своё комплексно-сопряжённое.}
\lectureframe{01:16:40}
\end{figure}

\begin{classroomqa}
\lecturetimestamp{01:19:19}
\audiencequestion{Является ли это квадратом косинуса угла между двумя состояниями?}
\lecturerresponse{Для вещественных единичных векторов в обычном евклидовом пространстве скалярное произведение равно косинусу угла. Здесь скалярное произведение, вообще говоря, комплексно, так что нет обычного пространственного угла, косинусом которого оно было бы. Полезным аналогом служит модуль перекрытия, а вероятность равна его квадрату модуля.}
\end{classroomqa}

\begin{classroomqa}
\lecturetimestamp{01:20:18}
\audiencequestion{Что, если приготовленное состояние — это другой собственный вектор \(M\) с другим собственным значением?}
\lecturerresponse{Тогда эти два собственных вектора ортогональны. Их перекрытие равно нулю, поэтому вероятность получить \(\lambda_A\) равна нулю. Приготовление состояния с одним определённым значением исключает другое определённое значение той же наблюдаемой.}
\end{classroomqa}

\begin{classroomqa}
\lecturetimestamp{01:21:23}
\audiencequestion{Может ли то же правило описывать приготовление, при котором ни один исход не является достоверным?}
\lecturerresponse{Да. Это общий случай: \(\ket B\) не обязано быть собственным состоянием \(M\). Его перекрытия с нормированными собственными векторами дают вероятности соответствующих исходов. Мы начинаем с примера с двумя исходами, прежде чем переходить к более сложным распределениям.}
\end{classroomqa}

В лекции это правило отрабатывается на нескольких приготовлениях, а не прячется внутри одной формулы.

Сначала приготовим \(\sigma_3=+1\):
\begin{equation}
  \ket B=\ket{u}=\begin{pmatrix}1\\0\end{pmatrix}.
\end{equation}
Измерим \(\sigma_1\). Для его собственного вектора с собственным значением \(+1\),
\begin{equation}
  \braket{x+}{u}
  =
  \frac{1}{\sqrt2}
  \begin{pmatrix}1&1\end{pmatrix}
  \begin{pmatrix}1\\0\end{pmatrix}
  =\frac{1}{\sqrt2},
\end{equation}
поэтому
\begin{equation}
  P(\sigma_1=+1\mid\sigma_3=+1)=\frac12.
\end{equation}
Замена \(\ket{x+}\) на \(\ket{x-}\) даёт ту же вероятность. Если поменять местами приготовление и измерение, тоже получится одна вторая, потому что квадрат модуля перекрытия симметричен.

Для действительно комплексного примера приготовим \(\ket{x+}\) и измерим \(\sigma_2\). Тогда
\begin{align}
  \braket{y+}{x+}
  &=
  \frac12
  \begin{pmatrix}1&-i\end{pmatrix}
  \begin{pmatrix}1\\1\end{pmatrix}
  =\frac{1-i}{2},
  \notag\\
  P(\sigma_2=+1\mid\sigma_1=+1)
  &=\left|\frac{1-i}{2}\right|^2
  =\frac14+\frac14
  =\frac12.
  \label{eq:l03-complex-probability-example}
\end{align}
Следовательно, если приготовить спин вдоль одной оси и измерять его вдоль перпендикулярной оси, то вероятности двух исходов равны.

\section{Указатели, состояния и произвольные оси}

Предыдущее предложение можно понять неправильно, если не уточнить терминологию. Магнитный момент имеет пространственное направление, которое мы будем называть \emph{указателем}. Кет — это вектор в комплексном пространстве состояний. Два пространственных указателя могут быть перпендикулярны, тогда как два вектора состояния могут быть ортогональны. Это разные утверждения о разных пространствах.

\begin{classroomqa}
\lecturetimestamp{01:29:25}
\audiencequestion{Если два физических направления спина образуют угол в девяносто градусов, должны ли их векторы состояния быть ортогональными?}
\lecturerresponse{Нет. Состояния со спином вверх вдоль перпендикулярных пространственных осей имеют квадрат перекрытия \(1/2\), а не нуль. Состояние, ортогональное состоянию со спином вверх вдоль одной оси, — это состояние со спином вниз вдоль той же оси. Термин \emph{перпендикулярный} мы сохраняем для указателей в обычном пространстве, а \emph{ортогональный} — для идеально различимых состояний.}
\end{classroomqa}

\begin{classroomqa}
\lecturetimestamp{01:31:25}
\audiencequestion{Могли бы мы использовать термин \emph{независимый} вместо \emph{ортогональный}?}
\lecturerresponse{Нет. Линейная независимость говорит лишь о том, что ни один вектор не является скалярным кратным остальных; она не требует нулевых скалярных произведений. Ортогональность — более сильное условие, и здесь она несёт физический смысл идеальной различимости.}
\end{classroomqa}

Пусть \(\hat n=(n_1,n_2,n_3)\) — единичный указатель в обычном пространстве:
\begin{equation}
  n_1^2+n_2^2+n_3^2=1.
\end{equation}
Оператор спина вдоль этого направления является матричнозначным аналогом скалярного произведения:
\begin{equation}
  \boldsymbol{\sigma}\cdot\hat n
  =
  n_1\sigma_1+n_2\sigma_2+n_3\sigma_3
  =
  \begin{pmatrix}
    n_3&n_1-in_2\\
    n_1+in_2&-n_3
  \end{pmatrix}.
  \label{eq:l03-spin-arbitrary-axis}
\end{equation}
Он эрмитов: диагональные элементы вещественны, а внедиагональные элементы комплексно-сопряжены друг другу.

\begin{figure}[htbp]
\centering
\includegraphics[width=0.84\linewidth]{lecture_03_figure_07.png}
\caption{Пространственный единичный указатель, операторная сумма и её явная эрмитова матрица вводятся одновременно.}
\lectureframe{01:38:00}
\end{figure}

Чтобы найти его возможные значения, сначала возведите его в квадрат, прежде чем подставлять матричные элементы. Порядок матриц существен, поэтому перекрёстные члены возникают как антикоммутаторы:
\begin{align}
  (\boldsymbol{\sigma}\cdot\hat n)^2
  &=
  n_1^2\sigma_1^2+n_2^2\sigma_2^2+n_3^2\sigma_3^2
  \notag\\
  &\quad
  +n_1n_2(\sigma_1\sigma_2+\sigma_2\sigma_1)
  \notag\\
  &\quad
  +n_2n_3(\sigma_2\sigma_3+\sigma_3\sigma_2)
  \notag\\
  &\quad
  +n_1n_3(\sigma_1\sigma_3+\sigma_3\sigma_1).
  \label{eq:l03-axis-square-expanded}
\end{align}
Различные матрицы Паули антикоммутируют:
\begin{equation}
  \sigma_i\sigma_j+\sigma_j\sigma_i=0
  \qquad(i\ne j).
  \label{eq:l03-pauli-anticommutation}
\end{equation}
В лекции один случай проверяется явно:
\begin{equation}
  \sigma_1\sigma_3
  =
  \begin{pmatrix}0&-1\\1&0\end{pmatrix},
  \qquad
  \sigma_3\sigma_1
  =
  \begin{pmatrix}0&1\\-1&0\end{pmatrix}
  =-\sigma_1\sigma_3.
  \label{eq:l03-order-check}
\end{equation}

\begin{figure}[htbp]
\centering
\includegraphics[width=0.82\linewidth]{lecture_03_figure_08.png}
\caption{Перемножение \(\sigma_1\) и \(\sigma_3\) в противоположных порядках даёт матрицы, различающиеся знаком минус.}
\lectureframe{01:44:50}
\end{figure}

Перекрёстные члены в уравнении~\eqref{eq:l03-axis-square-expanded} обращаются в нуль, тогда как \(\sigma_i^2=I\). Следовательно
\begin{equation}
  (\boldsymbol{\sigma}\cdot\hat n)^2
  =(n_1^2+n_2^2+n_3^2)I
  =I.
\end{equation}
Поэтому каждое собственное значение удовлетворяет \(\lambda^2=1\):
\begin{equation}
  \lambda=\pm1.
\end{equation}
Какую бы пространственную ось ни выбирал прибор, детектор всё равно имеет только два исхода \(+1\) и \(-1\). Связь между близкими направлениями проявляется не в виде промежуточного измеренного значения, а в вероятностях этих двух фиксированных ответов.
